
\chapter{INTRODUCTION}
Elliptic curve cryptography is a type of public-key encryption used in institutions and government computer system. The Diffie-Hellman key exchange enables the transfer of secure information on the ElGamal public key exchange. ElGamal digital signature ensures the authentication of information among individuals. To high-light the aforementioned, it is essential to understand attacks on discrete logs, this will help to maximize existing approaches to maintaining security. Using a recommended domain parameters of secp192k1, we generate a private key.
\section{Background of the study}
\cite{manoj2010cryptography} defined cryptography as the art and science of hiding messages in order to add secrecy to information security. The term ``Cryptography" was coined by combining two Greek words; ``Krypto," which means ``hidden," and ``Graphene," which means ``writing," which is where cryptography is thought to have begun. As civilizations advanced, tribes, clans, and kingdoms arose. This resulted in the emergence of power, disputes, supremacy, and politics. These ideas encouraged people's natural need to engage privately with a limited group of people, guaranteeing that cryptography progressed.\\
The beginning of cryptography can be traced back to the  Roman and Egyptian civilizations. Hieroglyph is the oldest cryptographic technique. The first documented hint of cryptography is the use of the term `Hieroglyph'. Around 4000 years ago, the Egyptians used hieroglyphic messages to communicate. This code was only known by the scribes who used to deliver communications on behalf of the rulers.\\\\
Years later, in 1976, Diffie and Hellman invented public key cryptography, which not only revitalized the field of cryptography, but also influenced the direction of computational number theory research. For the first time, the relative difficulty of various number theorem tasks became a reasonable concern.\\
The rivest-shamir-adleman (RSA) cryptosystem was the first feasible public key system, having been introduced in 1978. It spurred a flurry of research into factoring and primality testing, two related issues. A type of public key cryptography based on the discrete counterpart of the logarithm function spawned a second line of research in computational number theory.\\
discrete log cryptosystems (DLC) have not been around as long as RSA, but they have caught academics interest. The structure has prompted a lot of research due to the practical challenges that have arisen in discrete log cryptography.\\\\	
In \cite{Afreen2011} Neal Koblitz and Victor Miller  independently proposed the usefulness of elliptic curves in encryption in 1985. Since then, elliptic curve cryptography (ECC) has grown into a substantial branch of Public Key Encryption.\\\\
Public key cryptosystem (PKC) use public keys to encrypt data. Data is encrypted and decoded using different keys in a PKC system. Because one of the keys in a PKC system is made public and the key size determines the system's security.\\\\
Prior to PKC, mathematical problems were solved using prime factorization and discrete logarithm.\\\\
Elliptic curve cryptography has demonstrated that smaller key sizes can achieve the same level of security. For the most part, elliptic curve cryptography is concerned with the security of application systems, storage, computing performance, as well as CPU and GPU architecture.\\\\
Cryptography has been useful from the past and it is still now. Cryptography has undergone series of evolution that has led to cloud computing, data security issues, challenges in the cloud, DNA computing and DNA cryptography.\\
Cloud computing provide pervasive, easy, on-demand network connectivity, to a configurable pool of computational resources; of which include network infrastructures, servers, on \& off storage sites, applications  and services. All these can be accessed and released quickly and with low effort and  interaction with service providers or management effort.

\section{Problem Statement}
Elliptic curve cryptography has many benefits that have helped in the field of science and technology in terms of information secrecy and keeping of confidential data for governments and institutions. Nevertheless, it does have its negative aspects. Channel attacks and twist-security attacks are examples of such attacks. Both types try to compromise the ECC's private key security. This project seeks to address and highlight some possible solutions to these problems.\\\\

\section{Objectives of study}
This project's major goal is to better understand elliptic curves and how they relate to cryptography. Abstract algebra concepts such as groups and fields are examined first as a prerequisite for the thesis. \\ Other Objectives are;
\begin{itemize}
	\item To examine cryptographic applications of elliptic curves.
	\item To understand how private key is generated and the attacks on security.\\\\
\end{itemize}

\section{Significance of the study}
Elliptic curves can be found in a wide range of mathematical fields, from number theory to complex analysis. Elliptic curves can be used for encryption, digital signatures, pseudo-random generators, and a variety of other applications. They are also employed in a number of integer factorization methods with cryptographic applications, such as the Lenstra elliptic-curve factorization, illustrating the importance of elliptic curves in human lives in the 1980s and years after. This thesis would illustrate how to employ elliptic curves in various applications, particularly when they are defined over finite fields. It also examines how and why such curves are useful in the field of cryptography.

\section{Scope of the study}
This thesis is not intended to address a specific topic or to produce a novel finding; rather, as previously said, it intends to provide a grasp of elliptic curves, how they came to be, and how they are applied in elliptic curve cryptography. The subject of elliptic curves is a strong blend of abstract algebra, in particular group law and fields, as we will see in this paper.\\
At the moment, elliptic curves can be used for encryption, digital signatures, pseudo-random generators, and a variety of other applications. They are also employed in a number of integer factorization methods with cryptographic applications. Elliptic curves are employed in cryptographic protocols like Bitcoin and Austrian smart e-ID, to name a few examples. We will go over each of these applications. Elliptic curves are a common theme in algebraic geometry and cryptography, among other branches of mathematics. We can find number theory problems in the history of elliptic curves and complex theoretical functions.\\

\section{Organization of Study}  
This thesis is divided into five chapters, the first of which provides a broad summary of the entire thesis as well as the circumstances of the research. This chapter begins with an introduction, the study's history, a problem statement, and the study's goal. We outline the work connected to elliptic curve cryptography in the project's second chapter. The third chapter contains thorough information on elliptic curve cryptograph. Detailed results and analysis about the generation of private key cryptography are presented in chapter four. The conclusions and recommendations from the results and discussions have been outlined in chapter five.

